\section{Composition and Phase Stabilization}

Composition control is a cross-cutting facet because the nominal formula can change before, during, and after thermal processing. The collection includes title-level studies framed around Ga addition and cubic-phase stabilization \cite{shinawi2013stabilization}, Ga substitution and grain coarsening \cite{li2016ga}, and lithium volatilization with phase change during Al-containing LLZO processing \cite{montoya2024lithium}. These titles identify three different control objects---substitution, microstructural development, and firing-stage inventory---but they do not establish site occupancy, phase fraction, or a causal chain in the present review.

The defect question is not reducible to a fixed vacancy count. The stored abstract of a first-principles study calculates native defect populations across synthesis conditions and reports that donor substitution does not yield the same simple lithium-vacancy compensation under all modeled conditions \cite{squires2019native}. This is theoretical evidence about a conditional defect landscape, not direct experimental proof for every processed pellet. It nevertheless justifies retaining nominal composition, measured composition where available, and synthesis environment as separate comparison fields.

Experimental abstracts reinforce the need for that separation. A precursor-homogeneity study reports changes in compositional and structural evolution and connects spatial mixing of reactants and dopant species to reproducible material characteristics within its experiments \cite{parascos2022compositional}. A Nb-containing submicron-powder study reports phase, relative density, microstructure, total conductivity, and activation energy under stated firing conditions \cite{yan2020submicron}. An AlN-additive study on Ta-containing garnet jointly reports densification, transport, and electrochemical readouts and proposes reaction products as part of its interpretation \cite{zhang2022high}. In the present evidence hierarchy, the reported associations are usable; the detailed site chemistry or mediation mechanism would require full-text examination.

Zr-site and multi-element changes are similarly inseparable from processing. The local title for a dual-substitution study also names spark plasma sintering \cite{dong2019dual}. It would be invalid to attribute its outcome exclusively to composition or exclusively to densification without the full factorial context. The same restraint applies to title-only records on tetragonal-LLZO lithium stoichiometry \cite{ilina2014lithium} and lithium excess in cubic-LLZO synthesis \cite{norazman2021the}: they map M2 coverage but do not determine the sign or magnitude of an effect here.

For comparison, phase labels must therefore travel with at least four distinctions. First, nominal and measured compositions are not interchangeable. Second, a proposed substitution site is not treated as measured occupancy unless the checked source establishes it. Third, starting lithium excess is separated from lithium supplied or lost during firing. Fourth, a phase assemblage and its measurement method are retained rather than compressed into ``cubic.'' This record supports direct comparison only when the associated powder route, thermal history, environment, density, and conductivity definition are also compatible.
